% !TEX root = ../../main.tex

\section{Corpus Sintético Maybe}

El primer corpus sintético usa la mónada \texttt{Maybe}. Su objetivo es validar fenómenos estructurales del Renamer y del grafo de precedencia con programas pequeños, deterministas y fáciles de comparar. La descripción completa de los casos está en el archivo \texttt{cases.md} del corpus \texttt{maybe-monad}.

Las familias implementadas cubren los siguientes aspectos:

\begin{itemize}
    \item \texttt{dependency-shapes}: formas de grafo como cadenas, grafo vacío, diamante, fanout, fanin y cadenas independientes.
    \item \texttt{binders}: patrones que definen múltiples variables, patrones de lista, patrones anidados, wildcards, as-patterns y patrones lazy.
    \item \texttt{let-statements}: interacción entre statements \texttt{let}, binders monádicos y shadowing.
    \item \texttt{body-stmts}: statements sin binder, tanto independientes como dependientes.
    \item \texttt{maybe-failure}: comportamiento de \texttt{Nothing}, fallos de patrón y guards.
    \item \texttt{shadowing}: reutilización del mismo nombre fuente con nombres renombrados distintos.
    \item \texttt{cost-selection}: casos donde el orden original no es mínimo, todos los costos empatan o existe un mínimo único.
    \item \texttt{controls}: controles negativos y positivos para verificar que el marcador conmutativo activa la extensión solo cuando corresponde.
\end{itemize}

Este corpus permite aislar errores en la recolección de binders, el cálculo de variables libres locales, la construcción de aristas RAW y la selección de candidatos. El caso guía del informe pertenece a la familia \texttt{cost-selection}: su resumen reporta costo secuencial 4, costo con \texttt{ApplicativeDo} 3 y costo mínimo con reordenamiento 2.
